% GENERATED FILE. Edit canonical lessons, then run npm run generate:books.
% canonical-source-hash: fnv1a64:c834efef2813766b
% canonical-lessons: PT-C266-mais-do-que, PT-C266-ha-tres-anos-que, PT-C266-revisao-experiencia

\chapter{Mais Do Que --- More Than, and How Long}
\label{ch:pt-more-than-and-how-long}
\hlchaptermodality{266}

\begin{chapteropening}
\textbf{By the end of this chapter:} \emph{I can compare two things with mais … do que, and say how long something has been going on with há … que and desde.}

Complete PT-C266-revisao-experiencia: describe an experience by how it felt, what you hope, how it compared and how long it has lasted.
\end{chapteropening}

\section[mais … do que]{mais … do que — more than}

\label{lesson:PT-C266-mais-do-que}



\begin{quote}
Answer \emph{tudo bem?} with a shrug. (\emph{Mais ou menos.}) Literally, \textquotedblleft{}more or less\textquotedblright{}.

Those two words, \textbf{mais} and \textbf{menos}, are the whole of this lesson's new grammar.
\end{quote}

\subsection*{You'll want to know: mais … do que}
To compare two things, put \textbf{mais}, \textquotedblleft{}more\textquotedblright{}, in front of the adjective and \textbf{do que}, \textquotedblleft{}than\textquotedblright{}, in front of the second thing:

\begin{quote}
O comboio é \textbf{mais} rápido \textbf{do que} o autocarro. — The train is faster than the bus. O avião é \textbf{mais} caro \textbf{do que} o comboio. — The plane is dearer than the train. És \textbf{mais} alto \textbf{do que} eu. — You are taller than me.
\end{quote}

English often builds the comparison into the adjective, \emph{faster}, \emph{taller}; Portuguese nearly always builds it with \emph{mais}. After \emph{do que} a person stays in the plain form, \emph{eu}, \emph{tu}: \emph{do que eu}. You will also hear plain \emph{que}, \emph{mais rápido que o autocarro}, in Portugal and in Brazil alike.

Its opposite is \textbf{menos}, \textquotedblleft{}less\textquotedblright{}: \emph{O autocarro é \textbf{menos} caro \textbf{do que} o comboio} — the bus is less expensive than the train.

\begin{grammarlens}[title={Grammar Lens: o mais, and two that refuse}]
For the most of all, keep \emph{mais} and give the noun its article:

\begin{quote}
\emph{o comboio \textbf{mais} rápido} — the fastest train \emph{a cidade \textbf{mais} bonita} — the most beautiful city
\end{quote}

The adjective still comes after the noun, and the article does the work of English \emph{-est}. Who it is the most among comes with \emph{de}: \emph{a cidade mais bonita \textbf{de} Portugal}.

Two adjectives refuse \emph{mais}. \emph{Bom} becomes \textbf{melhor}, \textquotedblleft{}better\textquotedblright{}, and \emph{mau} becomes \textbf{pior}, \textquotedblleft{}worse\textquotedblright{}: \emph{o comboio é melhor do que o autocarro}. You have heard both already, inside \textbf{melhorar}, \textquotedblleft{}to get better\textquotedblright{}, and \textbf{piorar}, \textquotedblleft{}to get worse\textquotedblright{}.
\end{grammarlens}

\begin{cousinweb}
\textbf{mais} is inherited from Latin \emph{magis}, \textquotedblleft{}more\textquotedblright{} (EtymDB 2.1), and so is the \textbf{mas} of chapter two hundred and sixty: \textquotedblleft{}more\textquotedblright{} became \textquotedblleft{}rather\textquotedblright{}, and \textquotedblleft{}rather\textquotedblright{} became \textquotedblleft{}but\textquotedblright{}. \textbf{menos} has no parent in EtymDB 2.1; it is usually traced to Latin \emph{minus}, \textquotedblleft{}less\textquotedblright{} (Lewis \& Short), which English kept whole as \textbf{minus}. \textbf{melhor} is inherited from Latin \emph{melior}, \textquotedblleft{}better\textquotedblright{} (EtymDB 2.1).
\end{cousinweb}

\subsection*{Your turn}
Say these aloud:

\begin{itemize}
\raggedright
  \item \emph{O comboio é mais rápido do que o autocarro.}
  \item the other way round — \emph{O autocarro é menos rápido do que o comboio.}
  \item \emph{És mais alto do que eu.}
  \item \emph{o comboio mais rápido}, then \emph{a cidade mais bonita de Portugal}
\end{itemize}

\begin{tcolorbox}[breakable,colback=teal!4,colframe=teal!35!black,title={Before you move on}]
Say \textquotedblleft{}the plane is dearer than the train\textquotedblright{}. (\emph{\textbf{O avião é mais caro do que o comboio.}}) How do you say \textquotedblleft{}the fastest train\textquotedblright{}? (\emph{\textbf{O comboio mais rápido.}}) What does \emph{bom} become? (\emph{\textbf{Melhor.}})
\end{tcolorbox}
\section[há três anos que]{há três anos que — how long, and still going}

\label{lesson:PT-C266-ha-tres-anos-que}



\begin{quote}
Say \textquotedblleft{}the train is faster than the bus\textquotedblright{}. (\emph{O comboio é mais rápido do que o autocarro.}) Then \textquotedblleft{}three years ago\textquotedblright{}. (\emph{Há três anos.})

An experience has a length, and Portuguese measures it with that same \emph{há}.
\end{quote}

\subsection*{You'll want to know: há três anos que}
\begin{quote}
\textbf{Há três anos que estudo} português. — I have been studying Portuguese for three years. \textbf{Estudo} português \textbf{há três anos.} — The same, the other way round.
\end{quote}

Listen to the verb. English says \emph{have been studying}; Portuguese says \textbf{estudo}, the plain \textbf{present}, because the studying is still going on.

Now set it beside the \emph{há} you learned for \textquotedblleft{}ago\textquotedblright{}. With an event in the past, the same three words mean something else:

\begin{quote}
\emph{Fui a Lisboa há três anos.} — I went to Lisbon three years ago.
\end{quote}

So the verb decides: with the present, it is still going; with the perfeito, it happened that long ago.

\begin{grammarlens}[title={Grammar Lens: desde, and durante}]
\textbf{desde}, \textquotedblleft{}since\textquotedblright{}, which you have, names the starting point instead of the length, and it takes the present for the same reason:

\begin{quote}
\textbf{Moro} aqui \textbf{desde} janeiro. — I have lived here since January.
\end{quote}

When the stretch has \textbf{ended}, the verb goes into the perfeito and the length can take \textbf{durante}, \textquotedblleft{}for, during\textquotedblright{}:

\begin{quote}
\textbf{Morei} em Lisboa \textbf{durante} dois anos. — I lived in Lisbon for two years. \emph{(and then I left)}
\end{quote}

One English word, \emph{for}, splits in two. \emph{Há … que} with the present: still going on. \emph{Durante} with the perfeito: over. In Brazil you will also hear \emph{faz três anos que} for the first.
\end{grammarlens}

\begin{cousinweb}
\textbf{há} is a form of \textbf{haver}, inherited from Latin \emph{habēre}, \textquotedblleft{}to have\textquotedblright{} (EtymDB 2.1), so \emph{há três anos que estudo} is, word for word, \textquotedblleft{}it has three years that I study\textquotedblright{}. \textbf{desde} is a compound of \emph{des}, itself Latin \emph{dē} plus \emph{ex}, and \emph{de} (EtymDB 2.1): \textquotedblleft{}from, out of, from\textquotedblright{}. \textbf{durante} is inherited from Latin \emph{dūrāns}, \textquotedblleft{}lasting\textquotedblright{} (EtymDB 2.1), built on \emph{dūrus}, \textquotedblleft{}hard\textquotedblright{}: what lasts is what holds hard.
\end{cousinweb}

\subsection*{Your turn}
Say these aloud:

\begin{itemize}
\raggedright
  \item \emph{Há três anos que estudo português.}
  \item \emph{Moro aqui desde janeiro.}
  \item the stretch that is over — \emph{Morei em Lisboa durante dois anos.}
  \item the same \emph{há}, now \textquotedblleft{}ago\textquotedblright{} — \emph{Fui a Lisboa há três anos.}
\end{itemize}

\begin{tcolorbox}[breakable,colback=teal!4,colframe=teal!35!black,title={Before you move on}]
Which tense goes with \emph{há três anos que}, and why? (The \textbf{present} — it is still going on.) What does \emph{há três anos} mean with the perfeito? (\textbf{Three years ago.}) And which word takes over when the stretch is finished? (\emph{\textbf{Durante}}, with the perfeito.)
\end{tcolorbox}
\section[(review)]{Review — describe an experience}

\label{lesson:PT-C266-revisao-experiencia}



\begin{quote}
Say \textquotedblleft{}I have been studying Portuguese for three years\textquotedblright{}. (\emph{Há três anos que estudo português.})
\end{quote}

\begin{grammarlens}[title={What you've built}]
Two chapters gave you the four things that turn a fact into an \textbf{experiência}: how it felt, what you hope, how it compared, and how long it went on.

\begin{quote}
\emph{Há três anos que estudo português. Primeiro era difícil e eu estava nervoso, mas agora estou contente. Acho que o português é mais fácil do que a matemática. Espero falar bem.}
\end{quote}

Count the moves: a length still running, \emph{há três anos que}; a feeling in the past, \emph{estava nervoso}, against one now, \emph{estou contente}; a comparison, \emph{mais fácil do que a matemática}, marked as an opinion; a hope, \emph{espero falar}.
\end{grammarlens}

\subsection*{Your turn}
Say these aloud:

\begin{itemize}
\raggedright
  \item \emph{Há três anos que estudo português.}
  \item \emph{Primeiro era difícil e eu estava nervoso, mas agora estou contente.}
  \item \emph{Acho que o português é mais fácil do que a matemática.}
  \item a finished stretch — \emph{Morei em Lisboa durante um ano, e estava contente com a cidade.}
  \item your own experience, with all four moves, ending on \emph{O meu sonho é …}
\end{itemize}

\begin{tcolorbox}[breakable,colback=teal!4,colframe=teal!35!black,title={Before you move on}]
Name the four moves of an experience. (How it \textbf{felt}, what you \textbf{hope}, how it \textbf{compared}, how \textbf{long}.) Which tense goes with \emph{há … que}? (The \textbf{present}.) Say \textquotedblleft{}the most beautiful city\textquotedblright{}. (\emph{\textbf{A cidade mais bonita.}}) And when you tell the whole experience to someone, which verb counts it out? (\emph{\textbf{Contar}}, from Latin \emph{computāre}.)
\end{tcolorbox}
